% 2009 年辽宁卷文科第 20 题：2×2 列联表（待填）。
% 来源 PDF：普通高考/2009/2009辽宁文.pdf，SHA256 99623efa457632f26bc5d5775293bf9fd13a56add984179360e30698afc7cb73
% 原页：第 2 页（A3 三栏）第 2 栏中部，印刷页 540。
% 重绘依据：行标签为优质品、非优质品、合计，列标签为甲厂、乙厂、合计；表内单元格在原页留空，待填。
\begin{tikzpicture}[semithick, every node/.style={inner sep=2pt, font=\small}]
  \draw (0,0) rectangle (7.4,-2.8);
  \draw (2.6,0) -- (2.6,-2.8);
  \draw (4.2,0) -- (4.2,-2.8);
  \draw (5.8,0) -- (5.8,-2.8);
  \foreach \y in {-0.7,-1.4,-2.1} { \draw (0,\y) -- (7.4,\y); }
  \node at (1.3,-0.35) {};
  \node at (3.4,-0.35) {甲厂};
  \node at (5.0,-0.35) {乙厂};
  \node at (6.6,-0.35) {合计};
  \node at (1.3,-1.05) {优质品};
  \node at (1.3,-1.75) {非优质品};
  \node at (1.3,-2.45) {合计};
\end{tikzpicture}
